\section{Proof obligations}
\label{sec:obligations}

This section keeps two ledgers. The first, \cref{sec:obligations-current},
holds the obligations open against the current specification, CCME
3.1.0-rc.2 at repository commit \code{252216ba}, and against the reference
implementation at the same commit: places where the normative text
contradicts itself, where the text states a boundary of a guarantee rather
than a rule, and where the engine departs from the text. Each states the
claim, a minimal scenario, the consequence for the results of this report, the
repair or the recorded status, and the case and test that pin it. They are
ordered by severity, the silent loss of a release first and diagnostic
metadata last. The second ledger, \cref{sec:obligations-historical}, is the
audit of CCME 1.0.0 at \code{713f1a0b} that this report was first written for.
Its six obligations are provenance: the released 2.0.0 text incorporated
their repairs and the current text keeps them, so they are reproduced
unchanged.

Every \SPEC{} reference in \cref{sec:obligations-current} is to the current
text. In \cref{sec:obligations-historical}, except where a paragraph
explicitly names CCME 2.0.0 or a current theorem, every numbered reference
and quoted guarantee refers to the 1.0.0 snapshot, and present tense
describes the historical rule being executed.

\subsection{Current obligations}
\label{sec:obligations-current}

The sources consulted are \SPEC{} itself, its dated design history and
changelog, the release notes of the engine, the integration test plan, and
the test sources of the engine and of its integration suite. Where the design
history names a departure, the obligation quotes it; where it does not, the
obligation says what the code does and where.

\subsubsection{The owed window and \diag{E201} are specified but not implemented}

\begin{obligation}[a consumer that sits out the run in which its provider delivers is stranded]
\label{ob:owed}
\spec{13.4a} admits a bump-axis contribution for a target until the source
has \emph{delivered} it. \spec{13.3} adds, for every edge $D \to P$ where $D$
has a baseline, an owed window \code{owedWindow(P, D)} to the union of
windows parsed in \spec{13.4}, so that a debt survives the release of the
provider; and \spec{19.3} refuses with \diag{E201} a run that would release a
provider at the baseline commit of a consumer it still owes unless the
consumer releases after it in the same run. The reference implementation
implements the delivery test and the reconciliation of a proceeding consumer,
and neither the owed window nor \diag{E201}.
\end{obligation}

\paragraph{Source.} The design history entry of 22 September 2026,
``Delivery discharges a propagated contribution'', ends: ``The dispat engine
implements the delivery admission and the reconciliation of a proceeding
consumer; the owed windows and \code{E201} are not yet implemented, and a
consumer that sits out the run in which its provider releases the owed commit
is therefore still stranded in that engine, which its release notes list as a
departure.'' The engine's sources at the same commit contain no \code{E201}
and no owed window; \code{admission.go} carries \code{isDelivered} and
\code{owedSources}, the delivery half. The release notes at this commit do
not yet carry the departure: \code{services/dispat/CHANGELOG.md} ends at
\code{v1.11.0-rc.3} of 20 September, two days before the entry, and
\code{announce/rc/notes.yaml} does not mention it.

\paragraph{Scenario.} One commit $C_1$ carries \code{feat(core)\^{}: x} and
\code{feat(cli): y}, with \code{cli} consuming \code{core}. Run 1:
\code{core} fails to publish; \code{cli} has a cause of its own, proceeds at
\code{2.1.0} with its manifest naming \code{core}'s baseline \code{1.4.0},
and is tagged at $C_1$ (\spec{19.3}; vector 80d). A commit $C_2$ lands. Run 2
releases \code{core@1.5.0} at $C_2$ in a run \code{cli} sits out, held or
deselected. Now $W(\code{core}) = \emptyset$ and $W(\code{cli}) = \{C_2\}$:
$C_1$ lies in no window of either package, so the unit is never parsed,
\code{cli} is never planned, and its manifest names \code{core@1.4.0} for
ever. The delivery test would answer correctly if it were asked, since
\code{core}'s release commit $C_2$ is not reached by \code{cli}'s baseline
$C_1$; the owed window is what asks it. Vector 80d states the case in its
last sentence: ``\code{core} published at a later commit in a run \code{cli}
sat out leaves \code{C1} in \code{cli}'s owed window (\S13.3), so the next run
plans \code{cli}'s catch-up.''

\paragraph{Consequence.} In the specification, \cref{thm:g2} holds as
stated: the owed window keeps the contribution admissible until \code{cli}
releases at or after \code{core}'s release carrying $C_1$. In the engine, the
theorem holds only while some pending window still holds the commit, which
is every case but this one. The loss is silent in the plan. The tag-level
screen of \cref{prop:screen} does fire, because \code{core}'s stable commit
is not an ancestor of \code{cli}'s baseline, so a scheduled \spec{13.7b}
audit reports \diag{W195}; nothing releases it.

\paragraph{Repair.} Add the owed windows to the parsed union and \diag{E201}
to the publish protocol, as written. Until then the operator-side remedy is
the one \spec{19.3} states for the refused run: release the consumer in the
same run as the provider, or pin it with an explicit \code{Release-As}.

\paragraph{Pin.} Vector 80d, last two cases. The tests
\code{TestAdmissionCatchesUpAConsumerThatProceededOnItsOwn} and
\code{TestAdmissionCatchesUpAConsumerThatOvertookAHeldProvider} of the
integration suite and
\code{TestAdmissionCatchesUpAConsumerThatOvertookItsProvider} of
\code{internal/plan} pin the delivery half at an unchanged \Head{}; no test
in either suite exercises the run the consumer sits out.

\subsubsection{Reached-set attribution in the measured build}

\begin{obligation}[a unit's whole source set was credited to every reached dependent]
\label{ob:attribution}
\spec{9.2} attributes to a dependent $d$ the sources of a unit within the
unit's depth of $d$, and \spec{13.4a} draws the owed set from that reaching
set. The build that the campaigns of \cref{sec:evaluation} ran,
\code{1.11.0-rc.4+main.272772cb}, credited a unit's whole source set to every
dependent the walk reached: for a unit written over several packages, a
consumer two edges from one source and one from another was owed by both,
so it was released again on the far source's account after the near source
had delivered; and a consumer of only one of the packages was told it
releases because of the other. This is the (R) shape of
\cref{sec:guarantees} and the second shape of \cref{rem:owed-reach}.
\end{obligation}

\paragraph{Source.} Reported by the product session that ran the campaigns,
on 22 September 2026, as a correction of its own earlier statement that the
engine's attribution matched the specification; this report had not repeated
that statement. The fix is commit \code{e461d7d9} of the product repository's
\code{main}, ``attribute a unit's sources to a dependent within the unit's
depth of it'' (\code{internal/plan/admission.go}, \code{plan.go},
\code{propagate.go} and their tests; \code{owedSources} now answers a subset
of the reaching set in name order), after the snapshot \code{252216ba} this
report transcribes and unpublished when this was written. The commit message
records that \code{TestMultiScopeSourceOutsideTheManifestsStaysOutOfUpdates}
had pinned the old attribution and now expects the specified provenance.

\paragraph{Consequence for the evidence.} None for the figures of
\cref{sec:eval:rq5}: every trigger commit in the kernel and LLVM fixtures
scopes one package, and for a single-source unit the reaching set equals the
whole source set, so plans, tags and timings are identical under both rules,
and the reproduction kit's pin at \code{272772cb} stands. The companion
campaigns through 1.10.0 predate the delivery rule and are not affected.

\paragraph{Status.} Fixed after the measured build; the specification's text
still permits the shape, which \cref{rem:owed-reach} records as a
disagreement between its two formulations rather than an engine matter.

\subsubsection{Two releases on one commit: a stated boundary of completeness}

\begin{obligation}[the same-commit failure case of \diag{E201}]
\label{ob:samecommit}
\spec{19.3} states: ``If the consumer is published after the provider and
fails, the same state arises, and \code{E201} reports it with its remedy: an
explicit \code{Release-As: <version>} on the consumer (\S8.6), because no
later plan can compute the debt.'' The protocol reports this state; it does
not resolve it. Completeness (\cref{thm:g2}) therefore has a boundary the
specification itself states.
\end{obligation}

\paragraph{Scenario.} The state of \cref{ob:owed} after run 1, and run 2 at
the same \Head{} $= C_1$ with both packages in the plan: \code{core@1.5.0}
publishes and is tagged at $C_1$, then \code{cli@2.1.1} fails. Both
\code{core}'s release commit and \code{cli}'s baseline are $C_1$. By
\spec{13.4a}, \code{delivered(core, C1, cli)} asks whether some release
commit $t$ of \code{core} carrying $C_1$ lies in the reach of \code{cli}'s
baseline commit; $t = C_1$ does, so nothing is owed, although \code{cli}
still names \code{core@1.4.0}. \spec{13.4a} says why no rule over tags can do
better: ``Two releases on one commit have no order (\S13.7b), so
\code{delivered} cannot tell a target that released at commit \code{t} after
its source from one that released there in an earlier run whose source then
released at \code{t} too.''

\paragraph{Consequence.} \Cref{thm:g2} discharges a contribution when $d$
releases at a commit whose ancestry contains $\cof{u}$. Ancestry is a partial
order on commits and says nothing about two records on one commit, so the
theorem is exact only under the hypothesis that a provider is never released
at the baseline commit of a consumer it owes unless that consumer releases
after it in the same run and succeeds. \spec{19.3} enforces the first half
before publication and reports the second with a remedy. The closing remark
of \cref{prop:screen}, that this case is prevented rather than detected,
holds with the qualification that the prevention is a refusal and the
residue is a report.

\paragraph{Status.} A boundary, stated. In the engine it is unreported,
because \diag{E201} is not implemented (\cref{ob:owed}); the state arises
exactly as described and leaves no diagnostic. Vector 80d names the refused
run; no vector or test exercises the failure after the provider's
publication.

\subsubsection{The version summary still contradicts the graduation rule}

\begin{obligation}[\spec{13.9} does not preserve \spec{11.5}'s train target]
\label{ob:graduation}
\spec{11.5} states the graduation version as ``\code{applyBump(S, E)}, the
same \code{target} as \S11.4, with no prerelease suffix'', where \spec{11.4}
defines $E$ as ``the effective bump accumulated over the pending window
(\S13.3)'', the train window $W(P)$. \spec{13.9}'s stable branch instead
reads ``\code{next = applyBump(stableBaseline(P), effective(P))}'', and
\code{effective(P)} is \code{max(direct(P), propagated(P))} (\spec{13.7})
with \code{direct(P)} taken ``over surviving tuples for \code{P} whose commit
is in \code{Wfresh(P)}'' (\spec{13.6}). On a channel-only graduation the
fresh effective bump is \bnone{} although an older minor or major change
established the train; read literally, the summary loses the aggregate the
specific rule requires. This report follows \spec{11.5}; the summary should
delegate to it or carry a separately named train target.
\end{obligation}

\paragraph{What changed since 2.0.0.} The obligation as first recorded had a
second half: \spec{11.5} said \diag{E185} fires when the target is not
greater than the ``baseline core'', which would have refused every ordinary
graduation. The current text resolves it: ``Graduation never lowers a
version: if \code{target} is lower than the core of the baseline, \code{E185}
is raised. A \code{target} equal to that core is the ordinary graduation,
because \code{1.3.0} ranks above \code{1.3.0-rc.1}.'' The same section now
also carries the channel-entry patch on the way out of a train (\diag{W204}
on graduation, vector 39b), which removes the dead end in which a train
entered by the patch could not graduate. Both are implemented. The first
half of the obligation stands: \spec{13.6}'s own second paragraph says that
``\S11.4 then uses the full \code{W(P)} to compute that admitted
prerelease's aggregate target'', and the window table of \spec{13.3} assigns
the aggregate to $W(P)$, but \spec{13.9} is the section an implementer
transcribes and it names neither.

\paragraph{Consequence.} \Cref{thm:g3} proves that fixed version inputs give
the same result on retry; it does not decide between two formulas.
\Cref{def:plan} follows \spec{11.5}. A second implementation transcribing
\spec{13.9} alone graduates a channel-only train to its stable baseline,
which \spec{13.9}'s own last guard then refuses as \diag{E195}.

\paragraph{Pin.} The last row of \spec{11.4}'s table (``\code{1.3.0-rc.1};
nothing new, \code{E} is still the window's \code{minor}; \code{stable};
\code{1.3.0}; graduation'') and edge case 61 of \spec{15.5}: ``Graduation
with no pending bumps: publishes the accumulated \code{target}.'' The engine
follows the specific rule: \code{computeVersion} in \code{plan.go} applies
the bump held over the whole train to the stable baseline, and
\code{TestTrainGraduationCountsPublishedWork} asserts that ``the bump of the
already-published breaking change must still be in force even though it
alone would not have re-released the train'';
\code{TestGraduationUpdatesSpanTheTrain} and
\code{TestTrainEnteredByTheEntryPatchGraduates} pin the record and the
patch.

\subsubsection{One edge-kind set orders, blocks and refuses cycles}

\begin{obligation}[no \code{orderKinds} or \code{blockingKinds} in the engine]
\label{ob:kinds}
\spec{13.1} requires the graph to be acyclic ``over the edge kinds of
\code{propagation.kinds} and \code{publish.orderKinds}'', and adds that a
cycle ``existing only through [\code{devDependencies}] is not a cycle for
this purpose and does not trigger \code{E200}, which is the common and
legitimate case''. \spec{19.2} orders over \code{publish.orderKinds},
``which includes \code{optionalDependencies}''; \spec{19.3} blocks over
\code{publish.blockingKinds}, ``which excludes \code{optionalDependencies}''.
The engine has neither key. Every declared edge, whatever its kind, enters
the one graph that the cycle check, the publish order and the failure
closure read; only propagation honours a kind set.
\end{obligation}

\paragraph{Where.} \code{config.go} builds one \code{model.Dependency} per
declared edge, carrying its kind and dropping none; \code{loadWorkspace} in
\code{plan.go} adds each to the graph and reports \diag{E200} on any cycle
of it; the publish order is that graph's topological order; and the blocking
closure of \code{release/executor.go} reads a per-relation boolean,
\code{isBlocking}, not the edge's kind. The kind set of
\code{parser.propagation.kinds} is consulted by the propagation walk alone,
which is what the integration test plan states as the purpose of
\code{TestPropagationKindsSelectTheEdgesTraversed}. Neither the design
history, the changelog, the release notes nor the documentation records the
difference; the documentation of \code{propagation.kinds} describes the walk
only.

\paragraph{Scenario.} Two packages that list each other under
\code{devDependencies}, a test fixture depending back on the package it
exercises. The specification accepts the workspace (vector 100a, edge case
99a); the engine refuses it with \diag{E200} before any plan. Separately,
\code{docs-site} depending on \code{ui} under \code{optionalDependencies}
with \code{ui}'s publish failing: the specification publishes
\code{docs-site} (vector 79, edge case 100); the engine blocks it with
\diag{W194} unless its relation declares \code{isBlocking: false}.

\paragraph{Consequence.} The departure is on the safe side. The engine's
graph is a superset of the one \spec{19.2} orders, so any topological order
of it is an order the specification accepts and \cref{prop:order} is
unaffected; the blocked set of \cref{def:blocking} is computed over a
superset of \code{blockingKinds}, so it is a superset of the specified one
and the $F \cup B$ of \cref{thm:resumption} is larger, never smaller.
Nothing the specification forbids is published. What is lost is admission: a
workspace with a development-only cycle cannot release at all, and an
optional dependent waits for a provider it does not need.

\paragraph{Repair.} Either implement the two keys with the specified
defaults, or record the departure where the design history records the
others and restrict \diag{E200} to the runtime kinds. The second is a
one-line change to the graph the cycle check reads.

\paragraph{Pin.} Vectors 79 and 100a and edge cases 99a and 100 are the
conformance cases; none has a test, and both vectors fail against the engine
as it stands. \code{TestFatalDependencyCycle} of the integration suite and
\code{TestCycleIsRejected} of \code{internal/plan} pin \diag{E200} over
runtime edges; \code{TestExecutionDevDependencyOutputReachesConsumer} shows a
\code{devDependencies} edge ordering builds under \spec{28}, which is what
vector 22 of \spec{28.8} asks.

\subsubsection{The fleet recorder rejects a retry that annotates the same commit}

\begin{obligation}[the create-only rule's exception is implemented on one of two paths]
\label{ob:records}
\spec{19.1} states: ``A release tag is created, never replaced. Writing it
to the authoritative store MUST fail when the store already holds that tag
name, unless the stored tag names the same commit for the same package and
version, which is the retry of an uncertain write (\S19.4).'' The
single-repository recorder implements both halves; the fleet recorder of
\spec{27} implements the first and not the exception.
\end{obligation}

\paragraph{Where.} \code{gitx/records.go} leases every record on the name
being absent and, when the store refuses, reads the stored tag back peeled:
the same commit is ``this run's own write arriving twice and is a success'',
another commit is a published record left where it is (\diag{E221}). Its
inventory compares peeled lines because ``two runs annotating the same
commit write different tag objects, and a comparison on those would call one
release two''. The fleet recorder's \code{PushRelease} in
\code{gitx/polyrepo\_write.go} pushes each source tag with a plain unforced
\code{git push}, which Git rejects whenever the remote holds a different tag
object under that name, the same commit included; the rejection surfaces as
\diag{E335} after the publication.

\paragraph{Scenario.} One run's tag push is applied by the remote and its
acknowledgement is lost, or two checkouts of one source repository record
the same version at the same head in turn. The next push of that name
carries a new annotation of the same commit. In a single repository the push
is reported as \diag{W223}, ``tag already exists on the remote at this
release's commit, skipped''; in a fleet it is refused, the run reports
\diag{E335}, and the dependent legs are blocked (\spec{27.7}).

\paragraph{Consequence.} Nothing incorrect is recorded: the refused write is
the identity the store already holds, and the next run reads it as released.
The cost is a failed run and blocked consumers for a retry the specification
defines as successful. \Cref{cor:exactly-once} is unaffected, since the
ledger holds one record either way.

\paragraph{Status.} Open for the fleet path; resolved for the
single-repository path by the design history entry of 21 September 2026,
``Record authority under the lock''. Neither the design history, the
changelog nor the test plan records the fleet path's difference.
\code{TestReleaseRecordsCreateOnlyPush} pins the single-repository
behaviour (``the same commit is the retry of a lost write and succeeds'');
\code{TestReleaseRecordsInAComposedWorkspace} pins the fleet's comparison
under the lock, not its push. Cases: \spec{28.8} vectors 27 and 28; edge case
76a.

\subsubsection{The failure category travels only from a distributed lock loss}

\begin{obligation}[the class of \spec{28.9} is attached at the new sites only]
\label{ob:category}
\spec{28.9} requires implementations to ``expose stable machine-readable
categories, named by the identifier each row states''. The design history
entry of 22 September 2026 states the engine's position: ``the existing
\code{E220} to \code{E222}, \code{E335} and \code{E336} carry
\code{native-recording-or-lock} as their class; the field is attached where
a distributed run loses its lock, and not yet at the older recording and
lock sites, which is a departure to be closed by attaching it there.''
\end{obligation}

\paragraph{Consequence.} A reader that switches on the class, which is what
\spec{28.9} designs the class for, sees a lost lock in a distributed run and
not the same lost lock in a local one, nor a refused record. No release
outcome depends on it. \code{execution/codes.go} names the mapping in one
place, so the repair is to emit the field at the five older sites.

\paragraph{Pin.} Vector 14 of \spec{28.8} and
\code{TestExecutionLockLossStopsNewEffects} of the integration suite cover
the site that carries the field; row EXE-12 of the requirement matrix
asserts \diag{E226} with its \code{execution-authority} class. No test
asserts the class on a local \diag{E221} or \diag{E336}.

\subsubsection{Profile departures and inapplicable vectors}

\begin{obligation}[what the engine records as not implemented or done otherwise]
\label{ob:profile}
The design history entry of 22 September 2026 lists, beside
\cref{ob:category}, the remaining departures of the engine from \spec{28}
and the vectors that do not apply to it. They are recorded here so that the
table of \cref{sec:conformance} reads against them.
\end{obligation}

\begin{itemize}
  \item ``Not implemented, so the conditional vectors do not apply: reuse of
    verified outputs across runs (vector 20) and the rollback profile of
    section 26 (vector 26).'' \spec{28.8} states the rule the entry relies
    on: a vector whose condition does not hold for an implementation ``is
    neither satisfied nor violated''.
    \code{TestExecutionOldRunReceiptGrantsNothing} and
    \code{TestExecutionTaskReceiptDischargesNothing} pin the halves of those
    vectors that hold regardless.
  \item ``Departures from the profile as written: the platform check of
    section 28.2 runs before dispatch for the packages a run releases and,
    for a prepared provider, at its placement, where an unsatisfiable
    platform fails the provider's consumers.'' \spec{28.2} asks that
    unsupported target platforms ``fail preflight''.
    \code{TestExecutionPreparedProviderObeysItsOwnPlatforms} and
    \code{TestExecutionOrchestratorOnlyBuildIsNotHeldToTheWorkersPlatforms}
    pin the engine's timing.
  \item \spec{28.4} ``now also permits a durable mailbox ref or a
    history-free snapshot, and neither is implemented''. These are
    permissions, not requirements; the entry records the cost of the mode in
    use, a first transport of a whole history.
  \item The specification's own header keeps the VCS adapters of \spec{25}
    and the explicit rollback of \spec{26} ``pending implementation''. The
    forward-publish projection this report models is unaffected;
    \spec{13.7c} says its guarantees are claims about that projection.
\end{itemize}

\subsubsection{Textual defects noticed in passing}

Four defects of the text, none of which changes a rule, are recorded for the
specification's owner because they affect citation. Appendix B numbers two
vectors ``82c'': the consumer that gets ahead of a held provider (B.7) and
the hoisting of \code{resolvableBy} over a mixed source set (B.7a); this
report cites them as ``82c (B.7)'' and ``82c (B.7a)''. In \spec{19.3} the
bullets ``continue publishing packages that are not blocked'' and ``not
release a package at the baseline commit of a consumer it still owes'' are
run into the preceding bullet's line rather than set as items. The design
history says the release notes list the departure of \cref{ob:owed}, and at
this commit they do not. And the window table of \spec{13.3} cites
\spec{13.6} as a consumer of the aggregate train window, while the formula
of \spec{13.6} reads the undelivered one and only its prose defers the
aggregate to \spec{11.4}, the seam \cref{ob:graduation} sits on.
